% Options for packages loaded elsewhere
\PassOptionsToPackage{unicode}{hyperref}
\PassOptionsToPackage{hyphens}{url}
\PassOptionsToPackage{dvipsnames,svgnames,x11names}{xcolor}
\documentclass[
  11pt,
  a4paper,
]{article}
\usepackage{xcolor}
\usepackage[margin=2.6cm]{geometry}
\usepackage{amsmath,amssymb}
\setcounter{secnumdepth}{5}
\usepackage{iftex}
\ifPDFTeX
  \usepackage[T1]{fontenc}
  \usepackage[utf8]{inputenc}
  \usepackage{textcomp} % provide euro and other symbols
\else % if luatex or xetex
  \usepackage{unicode-math} % this also loads fontspec
  \defaultfontfeatures{Scale=MatchLowercase}
  \defaultfontfeatures[\rmfamily]{Ligatures=TeX,Scale=1}
\fi
\usepackage{lmodern}
\ifPDFTeX\else
  % xetex/luatex font selection
  \setmainfont[]{Libertinus Serif}
  \setmonofont[Scale=0.85]{Latin Modern Mono}
  \setmathfont[]{Latin Modern Math}
\fi
% Use upquote if available, for straight quotes in verbatim environments
\IfFileExists{upquote.sty}{\usepackage{upquote}}{}
\IfFileExists{microtype.sty}{% use microtype if available
  \usepackage[]{microtype}
  \UseMicrotypeSet[protrusion]{basicmath} % disable protrusion for tt fonts
}{}
\makeatletter
\@ifundefined{KOMAClassName}{% if non-KOMA class
  \IfFileExists{parskip.sty}{%
    \usepackage{parskip}
  }{% else
    \setlength{\parindent}{0pt}
    \setlength{\parskip}{6pt plus 2pt minus 1pt}}
}{% if KOMA class
  \KOMAoptions{parskip=half}}
\makeatother
\usepackage{longtable,booktabs,array}
\newcounter{none} % for unnumbered tables
\usepackage{calc} % for calculating minipage widths
% Correct order of tables after \paragraph or \subparagraph
\usepackage{etoolbox}
\makeatletter
\patchcmd\longtable{\par}{\if@noskipsec\mbox{}\fi\par}{}{}
\makeatother
% Allow footnotes in longtable head/foot
\IfFileExists{footnotehyper.sty}{\usepackage{footnotehyper}}{\usepackage{footnote}}
\makesavenoteenv{longtable}
\setlength{\emergencystretch}{3em} % prevent overfull lines
\providecommand{\tightlist}{%
  \setlength{\itemsep}{0pt}\setlength{\parskip}{0pt}}

\providecommand{\E}{\mathbb{E}}
\providecommand{\N}{\mathbb{N}}
\providecommand{\Z}{\mathbb{Z}}
\providecommand{\R}{\mathbb{R}}
\providecommand{\eps}{\varepsilon}
\providecommand{\ceil}[1]{\left\lceil #1\right\rceil}
\providecommand{\floor}[1]{\left\lfloor #1\right\rfloor}
\AtBeginDocument{\let\setminus\smallsetminus}
\usepackage[most]{tcolorbox}
\newtcolorbox{warning}{colback=red!5,colframe=red!65!black,boxrule=0.8pt,arc=2pt,left=6pt,right=6pt,top=5pt,bottom=5pt}
\usepackage{etoolbox}
\AtBeginEnvironment{longtable}{\small}
\setlength{\emergencystretch}{3em}
\usepackage{fancyhdr}
\pagestyle{fancy}
\fancyhf{}
\fancyhead[L]{\small\itshape Candidate result 1912.01570\_\_00 --- GPT-6 Astra, awaiting review}
\fancyhead[R]{\small\thepage}
\renewcommand{\headrulewidth}{0.2pt}
\usepackage{bookmark}
\IfFileExists{xurl.sty}{\usepackage{xurl}}{} % add URL line breaks if available
\urlstyle{same}
\hypersetup{
  pdftitle={A candidate counterexample to Conjecture 2 in Jones' Conjecture in subcubic graphs},
  pdfauthor={github.com/graph-theory-AI --- machine-generated candidate writeup},
  colorlinks=true,
  linkcolor={blue!50!black},
  filecolor={Maroon},
  citecolor={Blue},
  urlcolor={blue!50!black},
  pdfcreator={LaTeX via pandoc}}

\title{A candidate counterexample to Conjecture 2 in Jones' Conjecture
in subcubic graphs}
\author{\href{https://github.com/graph-theory-AI}{github.com/graph-theory-AI}
--- machine-generated candidate writeup}
\date{16 September 2026}

\begin{document}
\maketitle
\begin{abstract}
A 12-vertex planar graph has fvs = 5 and fp = 2; a related family has
fvs/fp tending to 3. This candidate disproof was generated by GPT-6
Astra in a single model pass for the Graph Theory AI project. It was
selected from the campaign because the model marked it
\texttt{would\_publish:\ true}; it has not been checked by the
adversarial referee, a human mathematician, or a novelty search.
\end{abstract}

\section{The problem and the claimed
result}\label{the-problem-and-the-claimed-result}

This document concerns
\href{https://graph-theory-ai.github.io/graph-conjectures/arxiv/1912.01570__00/}{Conjecture
2} from \emph{Jones' Conjecture in subcubic graphs}.

\subsection{Problem statement}\label{problem-statement}

For every planar graph \(G\), we have
\(\mathrm{fvs}(G) \leq 2 \cdot \mathrm{fp}(G)\), where
\(\mathrm{fp}(G)\) is the maximum size of a face-packing of \(G\), i.e.,
a cycle-packing where, for some embedding of \(G\), every cycle bounds a
face.

\subsection{Catalog context}\label{catalog-context}

The authors pose this stronger conjecture to highlight how little is
understood around Jones' Conjecture. They note that the example of many
nested disjoint cycles shows the embedding cannot be fixed, and that the
simple discharging argument giving
\(\mathrm{fvs}(G) \leq 3\,\mathrm{cp}(G)\) does not imply this
conjecture even with factor 3 instead of 2.

\subsection{Claimed result}\label{claimed-result}

\textbf{Model verdict: disproved} (confidence: high).

A 12-vertex planar graph has fvs = 5 and fp = 2; a related family has
fvs/fp tending to 3.

\emph{The model's own caveat:} Uses the stated vertex-disjoint,
planar-embedding definition; no claim of minimality or literature
priority.

\section{Model writeup}\label{model-writeup}

The text below is the model output verbatim, apart from moving its
machine-readable verdict block into the summary above and shifting
Markdown heading levels for this document. Mechanical TeX defects and
equation tags were normalized where needed for compilation.

\subsection{Counterexample theorem}\label{counterexample-theorem}

The conjecture is false as stated. More strongly, for every integer
\(t\ge 1\), there is a simple planar graph \(G_t\), of maximum degree
four and order \(8t+4\), such that \[
\mathrm{fvs}(G_t)=3t+2
\qquad\text{and}\qquad
\mathrm{fp}(G_t)=t+1.
\] Consequently, \[
\frac{\mathrm{fvs}(G_t)}{\mathrm{fp}(G_t)}
=3-\frac1{t+1}.
\]

For \(t=1\), this gives a 12-vertex counterexample: \[
\mathrm{fvs}(G_1)=5>4=2\,\mathrm{fp}(G_1).
\]

The important point is that the upper bound on face-packing below holds
\textbf{in every planar embedding}, not merely in a chosen drawing.

\subsection{1. Construction}\label{construction}

Take pairwise vertex-disjoint cycles \[
C_i=a_i b_i c_i d_i a_i,\qquad 1\le i\le t,
\] and pairwise vertex-disjoint copies \[
Q_j=G[\{p_j,q_j,r_j,s_j\}]\cong K_4,\qquad 0\le j\le t.
\] All these vertex sets are mutually disjoint.

Add exactly the following edges between them:

\begin{itemize}
\tightlist
\item
  \(a_i p_i\) and \(c_i q_i\), for \(1\le i\le t\);
\item
  \(d_i b_{i+1}\), for \(1\le i<t\);
\item
  \(b_1p_0\) and \(d_tq_0\).
\end{itemize}

Thus the \(C_i\)'s form a chain, each \(C_i\) has a \(K_4\) attached at
the opposite vertices \(a_i,c_i\), and \(Q_0\) connects the two ends of
the chain.

There are \[
4t+4(t+1)=8t+4
\] vertices. Every vertex has degree three except the vertices
\(p_j,q_j\), which have degree four.

For \(t=1\), the construction is particularly simple: take a square
\(abcd a\) and two disjoint \(K_4\)'s. Connect one \(K_4\) to \(a,c\),
and the other to \(b,d\), using distinct attachment vertices in each
\(K_4\).

\subsubsection{Planarity}\label{planarity}

First consider the auxiliary graph consisting of:

\begin{itemize}
\tightlist
\item
  the cycles \(C_i\);
\item
  the chain edges \(d_i b_{i+1}\);
\item
  a closing edge \(d_tb_1\);
\item
  the diagonals \(a_ic_i\).
\end{itemize}

Draw the \(C_i\)'s as diamonds in a row, with \(b_i,d_i\) their left and
right corners. Draw each diagonal \(a_ic_i\) inside its diamond and the
closing edge around the outside of the entire chain. This is a planar
drawing.

Now replace each diagonal \(a_ic_i\) by the two-terminal gadget \(Q_i\),
with connecting edges \(a_ip_i,c_iq_i\), and replace the closing edge by
\(Q_0\), with connecting edges \(b_1p_0,d_tq_0\).

These replacements preserve planarity: the two terminals \(p_j,q_j\) lie
on a common face of a plane \(K_4\), so each replacement can be
performed in a small neighborhood of the corresponding edge.

Choose the embedding of each \(Q_j\) with outer triangle \(p_jq_jr_j\)
and \(s_j\) inside. Its three internal triangular faces remain faces
after the replacement. In particular, \[
p_jq_js_jp_j,\qquad 0\le j\le t,
\] are pairwise vertex-disjoint facial cycles in this embedding.
Therefore \[
\mathrm{fp}(G_t)\ge t+1. \qquad\text{(1)}
\]

\subsection{2. Exact feedback vertex
number}\label{exact-feedback-vertex-number}

Every feedback vertex set must contain:

\begin{itemize}
\tightlist
\item
  at least one vertex of each \(C_i\);
\item
  at least two vertices of each \(Q_j\cong K_4\).
\end{itemize}

The latter assertion follows because any three vertices of a \(K_4\)
induce a triangle. These subgraphs have mutually disjoint vertex sets,
so \[
\mathrm{fvs}(G_t)\ge t+2(t+1)=3t+2. \qquad\text{(2)}
\]

For the reverse inequality, delete \[
S=\{a_i:1\le i\le t\}
  \,\cup\,
  \{p_j,q_j:0\le j\le t\}.
\] The remaining vertices of the \(C_i\)'s form the single path \[
b_1c_1d_1b_2c_2d_2\cdots b_tc_td_t,
\] and each \(Q_j\) leaves only the isolated edge \(r_js_j\). All
gadget-attachment edges have been removed. Thus \(G_t-S\) is a forest,
proving \[
\boxed{\mathrm{fvs}(G_t)=3t+2.} \qquad\text{(3)}
\]

\subsection{\texorpdfstring{3. A \(K_4\) gadget can be met by only one
packed
cycle}{3. A K\_4 gadget can be met by only one packed cycle}}\label{a-k_4-gadget-can-be-met-by-only-one-packed-cycle}

Fix \(j\). The only edges between \(Q_j\) and the rest of \(G_t\) are
its two attachment edges, incident with \(p_j\) and \(q_j\).

Consider any simple cycle meeting \(Q_j\).

\begin{itemize}
\tightlist
\item
  If it is contained in \(Q_j\), it uses at least three of its four
  vertices.
\item
  Otherwise, it must use both attachment edges, and hence contains both
  \(p_j\) and \(q_j\).
\end{itemize}

It follows that \textbf{any two cycles meeting \(Q_j\) intersect}:

\begin{itemize}
\tightlist
\item
  two cycles contained in \(Q_j\) intersect because each uses at least
  three vertices;
\item
  two cycles leaving \(Q_j\) both contain \(p_j,q_j\);
\item
  a cycle contained in \(Q_j\) intersects any cycle containing
  \(p_j,q_j\), since a three-element subset of a four-element set meets
  every two-element subset.
\end{itemize}

Therefore, in any vertex-disjoint cycle-packing, at most one cycle can
meet a given \(Q_j\). This observation is independent of the embedding.

\subsection{\texorpdfstring{4. None of the cycles \(C_i\) can ever be
facial}{4. None of the cycles C\_i can ever be facial}}\label{none-of-the-cycles-c_i-can-ever-be-facial}

Fix \(i\). There is an \(a_i\)--\(c_i\) path \[
P_i=a_ip_iq_ic_i
\] whose internal vertices lie in \(Q_i\).

There is also a \(b_i\)--\(d_i\) path \(R_i\) whose internal vertices
avoid both \(C_i\) and \(Q_i\): follow the chain from \(b_i\) towards
\(b_1\), cross through \[
b_1p_0q_0d_t,
\] and follow the other end of the chain back to \(d_i\). When
traversing another \(C_j\), use its path \(b_jc_jd_j\).

The paths \(P_i\) and \(R_i\) are internally vertex-disjoint, and their
endpoints alternate around \[
C_i=a_i b_i c_i d_i a_i.
\]

In any planar embedding, \(C_i\) is a Jordan curve. Two internally
disjoint paths with alternating endpoints on this curve cannot both lie
on the same side: a path joining \(a_i\) to \(c_i\) separates that side
into regions incident with \(b_i\) and \(d_i\), respectively.

Hence \(P_i\) and \(R_i\) lie on opposite sides of \(C_i\) in
\textbf{every} planar embedding. Both sides therefore contain graph
edges, so \[
\boxed{C_i\text{ is not facial in any planar embedding of }G_t.} \qquad\text{(4)}
\]

Equivalently, \(C_i\cup P_i\cup R_i\) is a subdivision of \(K_4\), with
\(C_i\) corresponding to a four-cycle rather than a facial triangle.

\subsection{5. Exact face-packing
number}\label{exact-face-packing-number}

Delete the vertices of all gadgets \(Q_j\). What remains is the chain of
cycles \(C_1,\ldots,C_t\), joined by the edges \(d_i b_{i+1}\). Those
joining edges are bridges in the remaining graph. Consequently, its only
simple cycles are the \(C_i\)'s.

By (4), none of those cycles can be facial in any planar embedding of
the full graph. Thus \textbf{every facial cycle of \(G_t\) must meet at
least one gadget \(Q_j\)}.

By Section 3, different members of a vertex-disjoint face-packing cannot
meet the same gadget. There are only \(t+1\) gadgets, so every planar
embedding satisfies \[
\text{maximum face-packing size}\le t+1.
\] Together with the explicit facial packing in (1), this proves \[
\boxed{\mathrm{fp}(G_t)=t+1.} \qquad\text{(5)}
\]

\subsection{6. Consequences and scope}\label{consequences-and-scope}

Taking \(t=1\) proves the claimed disproof: \[
\mathrm{fvs}(G_1)=5,
\qquad
\mathrm{fp}(G_1)=2.
\]

Moreover, the family shows that \textbf{no universal inequality} \[
\mathrm{fvs}(G)\le c\,\mathrm{fp}(G)
\] can hold for all planar graphs with any constant \(c<3\). This does
not establish an upper bound with constant three; it only shows that
such a constant would be best possible if an upper bound exists.

For comparison, \[
\mathrm{cp}(G_t)=2t+1.
\] The lower bound comes from taking all \(C_i\)'s and one triangle from
each \(Q_j\). The upper bound follows from Section 3 and the fact that
the only cycles avoiding all gadgets are the \(C_i\)'s. Thus these
examples do \textbf{not} contradict the original cycle-packing version
of Jones' conjecture.

The argument is self-contained and handles all planar embeddings. I have
not checked whether this construction has previously appeared in the
literature.

\appendix

\section{Provenance and review
status}\label{provenance-and-review-status}

\begin{warning}

\textbf{UNREFEREED MODEL OUTPUT.} This document was selected solely
because GPT-6 Astra labelled its own result
\texttt{would\_publish:\ true} and returned \texttt{proved} or
\texttt{disproved}. It has not passed the adversarial LLM referee used
for the earlier Sol campaign, has not been checked by a human
mathematician, and has not been checked for novelty. Treat every
mathematical and bibliographic claim as unverified.

\end{warning}

{\def\LTcaptype{none} % do not increment counter
\begin{longtable}[]{@{}
  >{\raggedright\arraybackslash}p{(\linewidth - 2\tabcolsep) * \real{0.5000}}
  >{\raggedright\arraybackslash}p{(\linewidth - 2\tabcolsep) * \real{0.5000}}@{}}
\toprule\noalign{}
\endhead
\bottomrule\noalign{}
\endlastfoot
Catalog id & \texttt{1912.01570\_\_00} \\
Catalog entry &
\href{https://graph-theory-ai.github.io/graph-conjectures/arxiv/1912.01570__00/}{Conjecture
2} \\
Source corpus & arxiv \\
Campaign leg & previously unattacked arXiv records
(\texttt{attacks\_arxiv\_astra}) \\
Source paper / entry & Jones' Conjecture in subcubic graphs \\
Model verdict & \textbf{disproved} (confidence: high) \\
Selection criterion & \texttt{would\_publish:\ true} and verdict
\texttt{proved} or \texttt{disproved} \\
Model & \texttt{gpt-6-astra}, reasoning effort \texttt{max},
\texttt{mode=pro}, flex service tier \\
Original artifact &
\texttt{attacks\_arxiv\_astra/1912.01570\_\_00/output.md} \\
Independent review & \textbf{None} \\
arXiv & \href{https://arxiv.org/abs/1912.01570}{arXiv:1912.01570} \\
Model's caveats & Uses the stated vertex-disjoint, planar-embedding
definition; no claim of minimality or literature priority. \\
\end{longtable}
}

The generated Markdown and LaTeX sources for this PDF are stored under
\texttt{to\_review\_astra/src/1912.01570\_\_00/}.

\end{document}
